\begin{abstract}


Commercial LLM agents retain user-provided information across conversations and offer various data deletion methods, yet it remains unclear whether these methods achieve the outcomes users expect.
%
In this study, we investigate users' data deletion practices, expectations, and verification strategies, alongside the post-deletion behavior of commercial LLM agents.
%
We conducted a survey of 177 U.S. adults with prior data deletion experience and developed a Data Injection-Deletion-Recall (DIDR) pipeline to evaluate whether synthetic user information remains retrievable after deletion.
%
Using DIDR, we conducted 865 test cases across six commercial LLM agents, with recall tested at least 30 days after deletion.
%
Our survey reveals that users generally expect deletion methods to permanently remove their information and prevent future retrieval, yet many lack reliable ways to verify deletion outcomes.
%
Our experiments further reveal substantial differences in post-deletion behavior across agents, with information sometimes remaining retrievable despite deletion requests and confirmations.
%
Moreover, post-deletion outcomes varied more across agents than across deletion methods or request wording.
%
By comparing users' expectations with observed agent behavior, we identify discrepancies between expected deletion outcomes and subsequent information retrievability, while finding that agents' responses and interface changes do not reliably indicate recall outcomes.
%
Our findings motivate clearer communication of deletion scope, verifiable deletion feedback, and standardized evaluation protocols to improve data deletion practices in LLM agents.


\end{abstract}
